% !TeX root = ../tesis.tex

%-----------------------------------------------------------------------
\section{Planteamiento del problema}
\label{sec:planteamiento}
%-----------------------------------------------------------------------

La problemática expuesta anteriormente revela que la Zona Metropolitana del
Valle de México (ZMVM) padece una desarticulación topográfica profunda, donde el
sistema de transporte masivo estructurado presenta una cobertura insuficiente en
las franjas periféricas que componen el \textbf{Tercer y Cuarto Contorno}
metropolitano \citep{semovi2020diagnostico}. Esta limitación técnica genera una
dependencia crítica del transporte público concesionado de baja capacidad y del
automóvil particular, lo cual resulta ineficiente para mitigar una demanda de
movilidad que desborda las fronteras administrativas de la capital.

Diversos diagnósticos urbanos coinciden en que la red masiva (Metro y Metrobús)
se ha consolidado bajo una lógica de centralización, obligando a los habitantes
de zonas periféricas en el \textbf{Norte (Ecatepec, Naucalpan)},
\textbf{Oriente (Iztapalapa, Chalco)} y \textbf{Sur (Tlalpan, Xochimilco)} a
realizar desplazamientos radiales que saturan los nodos del núcleo central
\citep{obras2025}. Esta configuración no solo incrementa los tiempos de traslado
---superando en promedio los 122 minutos para los habitantes del cuarto
contorno---, sino que eleva la \textit{vulnerabilidad funcional} del sistema
ante fallas operativas en hubs críticos como Pantitlán, Indios Verdes o Taxqueña
\citep{loterovélez2014}.

Aunque el Anillo Periférico fue proyectado como una infraestructura estratégica
para conectar la capital, su diseño priorizó históricamente la automovilidad
privada, dejando una brecha topográfica para el transporte masivo transversal.
La adopción de geometrías anillares ha demostrado ser una alternativa robusta en
grandes metrópolis para descomprimir el centro y articular las periferias
\citep{efe2023}; de este modo, surge la necesidad de contar con herramientas que
permitan evaluar cuantitativamente el impacto topológico de escenarios anillares
hipotéticos en la red metropolitana.

%-----------------------------------------------------------------------
\subsection{La morfología y topología en la periferia de la ZMVM}
\label{subsec:morfologia}
%-----------------------------------------------------------------------

La estructura urbana de la periferia metropolitana no constituye una extensión
continua y ordenada de la traza central. Siguiendo el análisis de
\citet{ortiz2007}, la morfología de la Ciudad de México se define como una
jerarquía de centros que da lugar a un ``\textit{Patchwork of offset grids}''
(mosaico de retículas desfasadas). En este modelo, cada sector periférico
(Norte, Oriente, Poniente o Sur) actúa como un ``parche'' geométricamente
aislado, cuya trama interna no se alinea con la estructura vial primaria de la
ciudad.

Esta fragmentación morfológica es el factor determinante de la ineficiencia en
las alcaldías y municipios del cinturón metropolitano. A diferencia de la
\textbf{Ciudad Central}, donde la conectividad en red facilita desplazamientos
multidireccionales, estos parches periféricos operan como islas conectadas
únicamente por arterias radiales que tienden al colapso crónico. Desde la
perspectiva del motor \texttt{VFTModel}, esta desconexión implica que los
caminos mínimos entre dos nodos periféricos aledaños tienen un
\textbf{Índice de Ruta Directa ($DI$)} significativamente elevado, ya que la red
vial obliga a los usuarios a converger hacia el centro para luego retornar a su
zona de destino \citep{fortin2016innovative}.

Complementariamente, la teoría de redes de \citet{owais2020} sostiene que en
ciudades con estructuras no ortogonales, la carencia de conexiones transversales
entre estos parches fragmentados genera una pérdida de robustez sistémica. Esta
teoría valida la tesis central de este trabajo: la problemática de la periferia
de la ZMVM no es únicamente una cuestión de capacidad vehicular, sino una
\textbf{desconexión geométrica estructural}. Mediante el procesamiento de datos
en \texttt{Apimetro}, se demostrará que la implementación de una ruta anillar en
el Tercer Contorno permite ``coser'' estos parches urbanos, optimizando la
topología del grafo metropolitano y reduciendo la fricción vial global.

